% ID: SEL-D0060-OB-E01
% ESTADO: validada. S060.
\textbf{a)} Las dos últimas restricciones equivalen a $x/4\le y\le3x/2$. Esto fuerza $x\ge0$ y ambas rectas se cruzan en el origen. Al cortar la frontera $13x+8y=600$ con $y=3x/2$ sale $(24,36)$, y con $y=x/4$ sale $(40,10)$. Los tres vértices satisfacen todas las restricciones.
\begin{center}
\setlength{\unitlength}{0.9mm}
\begin{picture}(94,56)(-7,-8)
\qbezier(0,0)(39,0)(78,0)
\qbezier(0,0)(0,23)(0,46)
\qbezier(0,0)(18,18)(36,36)
\qbezier(36,36)(48,23)(60,10)
\qbezier(60,10)(30,5)(0,0)
\put(79,0){$x$}
\put(2,45){$y$}
\put(35,38){\scriptsize $(24,36)$}
\put(62,10){\scriptsize $(40,10)$}
\put(-6,-6){\scriptsize $(0,0)$}
\end{picture}
\end{center}
\textbf{b)} Como $F=5(13x+8y)$, se tiene $F\le3000$. Hay igualdad precisamente en los puntos factibles de la recta $13x+8y=600$. El máximo es $\boxed{3000}$ en todo el segmento entre $(24,36)$ y $(40,10)$.
